%% CAPITULO 4 -- ANALISE E DISCUSSAO (esqueleto da dissertação)

\section{Análise e discussão}

A análise testará primeiro se estratégias equivalentes nas métricas tradicionais
de recuperação e de geração se diferenciam no IRD, o que sustentaria a hipótese
de que essas métricas não discriminam a adequação institucional da resposta. Em
seguida, examinará se o ganho das estratégias estruturadas depende do tipo de
pergunta, se o custo em \textit{tokens} cresce mais do que a qualidade entre
condições adjacentes e se os metadados melhoram a completude das respostas a
perguntas que dependem da posição da informação no documento.

Os resultados alimentarão a matriz de decisão do produto técnico, que indicará
em que cenários uma estratégia simples basta e em quais uma estratégia
estruturada se justifica pelo ganho de rastreabilidade. Como a pesquisa não
conta com implantação em um órgão, a matriz será avaliada quanto à eficácia, em
experimento contra uma linha de base de métrica única, e de forma formativa com
dois representantes dos perfis usuários, sem alegação de efetividade em uso.
